% Ampliación acumulativa del TRIT; se inserta antes de la figura de regímenes.
% No sustituye ni elimina el desarrollo precedente.
\subsubsection{División balanceada y composición de los levantamientos}
\label{trit-ampl-div-balanceada}

La reducción ternaria permite expresar una evaluación entera de APP en una
coordenada visible y una coordenada de prolongación. Su función no consiste
únicamente en abreviar la escritura de un entero: proporciona una regla de
composición que determina cómo se modifica el cociente cuando se suman o se
multiplican evaluaciones. La elección balanceada hace explícita, además, la
compatibilidad de esa regla con la inversión de signo.

\begin{definition}[Coordenadas ternarias balanceadas]
\label{trit-ampl-def-coordenadas}
Para \(n\in\mathbb Z\), sean
\[
 q_3(n)=\left\lfloor\frac{n+1}{3}\right\rfloor,\qquad
 r_3(n)=n-3q_3(n).
\]
La aplicación
\[
 \mathcal B:\mathbb Z\longrightarrow\TRIT\times\mathbb Z,\qquad
 n\longmapsto(r_3(n),q_3(n))
\]
se denomina descomposición ternaria balanceada. Su inversa es
\(\mathcal L(r,c)=r+3c\).
\end{definition}

En efecto, \(r_3(n)\in\{-1,0,+1\}\). Si dos pares representan el mismo entero,
\(r+3c=s+3d\), entonces \(r-s\) es múltiplo de \(3\) y pertenece al intervalo
\([-2,2]\); necesariamente \(r=s\) y \(c=d\). La representación es, por tanto,
única. La división puede repetirse sobre \(q_3(n)\), de manera que cada nivel
conserva el dato necesario para reconstruir el anterior. El cociente no es una
perturbación del residuo: constituye su coordenada complementaria.

La firma obtenida de una evaluación \(F\), con \(F=S\) o \(F=P\) en APP,
es \(r_3\circ F\). Cuando la reducción actúa sobre el cursor de fase, la misma
sección balanceada produce las clases \(\{1,4,7\}\), \(\{2,5,8\}\) y
\(\{3,6,9\}\), con valores \(+1,-1,0\). El selector del TPK que se define
a continuación utiliza esas clases para determinar el modo operativo.
La lectura de una evaluación y la selección de un modo tienen así la misma
reducción ternaria, pero dominios y funciones diferentes.

\Needspace{16\baselineskip}
\begin{proposition}[Aritmética de los pares balanceados]
\label{trit-ampl-prop-aritmetica}
Sean \(n=r+3c\) y \(m=s+3d\), con \(r,s\in\TRIT\). Defínase
\[
 \chi(r,s)=\frac{r+s-r_3(r+s)}3.
\]
Las operaciones enteras se transportan a los pares mediante
\begin{align}
 (r,c)\boxplus(s,d)
 &=\bigl(r_3(r+s),c+d+\chi(r,s)\bigr),
 \label{trit-ampl-eq-suma}\\
 (r,c)\boxtimes(s,d)
 &=\bigl(rs,rd+sc+3cd\bigr).
 \label{trit-ampl-eq-producto}
\end{align}
La aplicación \(\mathcal L\) conserva ambas operaciones.
\end{proposition}
\begin{proof}
La igualdad \(r+s=r_3(r+s)+3\chi(r,s)\) da
\[
 n+m=r_3(r+s)+3\bigl(c+d+\chi(r,s)\bigr).
\]
Por otra parte,
\[
 nm=rs+3(rd+sc+3cd).
\]
Como el producto de dos representantes balanceados pertenece de nuevo a
\(\{-1,0,+1\}\), no necesita una reducción residual adicional. La unicidad
de la descomposición demuestra las dos fórmulas.
\end{proof}

Por ejemplo, \(2=(-1)+3\) y \(4=1+3\). Su suma y su producto se reconstruyen
como
\[
 (-1,1)\boxplus(1,1)=(0,2),\qquad
 (-1,1)\boxtimes(1,1)=(-1,3).
\]
Las evaluaciones resultantes son \(6\) y \(8\). La suma de cocientes basta
en este ejemplo aditivo; en el producto intervienen los términos cruzados
\(rd\), \(sc\) y \(3cd\). La discrepancia entre ambas operaciones reside
también en la coordenada de prolongación, no solamente en el residuo.

El término \(\chi\) es un cociclo de la sección balanceada. Si
\(r\oplus s=r_3(r+s)\), satisface
\[
 \chi(r,s)+\chi(r\oplus s,t)
 =\chi(s,t)+\chi(r,s\oplus t).
\]
Ambos miembros representan el cociente de la misma suma de tres enteros.
Esta identidad garantiza que reagrupar las operaciones no cambia la
evaluación reconstruida. Al mismo tiempo, la proyección sobre el primer
componente suprime esa contabilidad: \(1+1=-1+3\), mientras que la lectura
residual conserva únicamente \(-1\). La información eliminada está
determinada, en este caso, por \(\chi(1,1)=1\).

La relación con el módulo \(9\) conserva una segunda escala discreta.
La aplicación
\[
 \beta:\TRIT\times\mathbb F_3\longrightarrow\mathbb Z/9\mathbb Z,
 \qquad (r,\bar c)\longmapsto r+3c\pmod9
\]
está bien definida y es biyectiva. Cambiar el representante de \(\bar c\)
añade un múltiplo de \(9\); recíprocamente, una igualdad de imágenes fija
primero \(r\) por reducción módulo \(3\) y después \(\bar c\). La suma
transportada por esta biyección contiene el término
\(\chi(r,s)\) de \eqref{trit-ampl-eq-suma}, reducido ahora módulo \(3\).
Así, los dos niveles ternarios no se componen mediante una suma
independiente de coordenadas: la normalización del primer nivel modifica
el segundo.

En particular, las clases \(3\), \(6\) y \(0\) de
\(\mathbb Z/9\mathbb Z\) tienen primer componente \(0\), pero segundo
componente \(1\), \(2\) y \(0\), respectivamente. La distinción ya observada
entre reducción directa y extracción de la coordenada del ideal aparece
como una propiedad de \(\beta\). Conservar esa coordenada impide que el
sector residual nulo se convierta artificialmente en una sola posición.
El módulo \(1000\), utilizado después para organizar bloques decimales,
cumple otra función de representación posicional; su cociente se conserva
mediante la división correspondiente. Un cambio de base conserva el entero
cuando retiene residuo y cociente, mientras que la coincidencia de residuos
en bases distintas no identifica por sí sola sus estados de transporte.

\Needspace{11\baselineskip}
\subsubsection{Inversión, clase residual y coordenadas no visibles}
\label{trit-ampl-inversion}

La inversión aditiva se levanta sin ambigüedad:
\[
 \mathcal B(-n)=(-r_3(n),-q_3(n)).
\]
En particular, \(3\), \(0\) y \(-3\) tienen coordenadas
\((0,1)\), \((0,0)\) y \((0,-1)\). Comparten firma nula, pero representan
evaluaciones distintas. En una transición restringida a un cociente de
amplitud máxima \(1\), son las tres posibilidades de la fibra visible sobre
cero; en un levantamiento entero general, esa fibra contiene todos los pares
\((0,c)\), con \(c\in\mathbb Z\).

La coordenada de cociente recupera el entero evaluado, pero no recupera por
sí sola la trayectoria que lo produjo. Dos caminos de APP pueden compartir
evaluación, residuo y cociente, y diferir en orientación, sucesión de modos o
enrollamiento. Por ello la descomposición \(\mathcal B\) se integra en el
estado de transporte, en vez de sustituirlo. La memoria aritmética de una
división y la memoria de una trayectoria son coordenadas relacionadas, no
sinónimos.

Para una sucesión de firmas de longitud \(n\), la inversión orientada es
\[
 C_{\rm rev}(\tau_1,\ldots,\tau_n)
   =(-\tau_n,\ldots,-\tau_1).
\]
Aplicarla dos veces restituye la sucesión. La operación invierte el orden de
las posiciones y el signo de cada firma; conserva su longitud y las
posiciones nulas con el orden invertido. En la parametrización adoptada,
intercambia el retorno con memoria, asociado a \(+1\), y la apertura
bidireccional, asociada a \(-1\), manteniendo el umbral \(0\). El transporte
de los restantes datos de una trayectoria se especifica mediante la
involución correspondiente del TPK.

Esta inversión de firmas debe distinguirse de la conjugación dentro de un
álgebra cuadrática. La primera puede cambiar el tipo \(\tau\); la segunda,
que se construye a continuación, actúa dentro del tipo ya fijado. La
distinción impide atribuir a una inversión de signo una equivalencia
algebraica entre regímenes diferentes.

\subsubsection{Producto cuadrático, conjugación y norma algebraica}
\label{trit-ampl-norma}

Un elemento de \(\mathbb A_\tau\) tiene expresión única \(aI+bJ_\tau\).
La relación \(J_\tau^2=-\tau I\) determina su producto:
\begin{equation}
 (aI+bJ_\tau)(cI+dJ_\tau)
   =(ac-\tau bd)I+(ad+bc)J_\tau.
 \label{trit-ampl-producto-cuadratico}
\end{equation}
Se trata de una realización real del tipo ya seleccionado. Las coordenadas
\(a,b\) no son residuos de APP, y el producto
\eqref{trit-ampl-producto-cuadratico} no sustituye a
\eqref{trit-ampl-eq-producto}; expresa otra operación, en el codominio
cuadrático declarado.

\begin{definition}[Conjugación y norma algebraica]
\label{trit-ampl-def-norma}
La conjugación del tipo \(\tau\) y su norma algebraica son
\[
 \overline{aI+bJ_\tau}=aI-bJ_\tau,\qquad
 N_\tau(aI+bJ_\tau)=a^2+\tau b^2.
\]
\end{definition}

\begin{proposition}[Multiplicatividad e invertibilidad]
\label{trit-ampl-prop-norma}
Para \(z,w\in\mathbb A_\tau\),
\[
 z\bar z=N_\tau(z)I,\qquad
 N_\tau(zw)=N_\tau(z)N_\tau(w).
\]
Además, \(z\) es invertible si y sólo si \(N_\tau(z)\ne0\), y en ese caso
\[
 z^{-1}=\frac{\bar z}{N_\tau(z)}.
\]
\end{proposition}
\begin{proof}
El producto de \(aI+bJ_\tau\) y \(aI-bJ_\tau\) es
\((a^2+\tau b^2)I\). La conjugación es multiplicativa y el álgebra es
conmutativa, por lo que
\((zw)\overline{zw}=(z\bar z)(w\bar w)\). Si la norma no se anula, la
fórmula indicada proporciona la inversa. Recíprocamente, si \(zw=I\), la
multiplicatividad implica \(N_\tau(z)N_\tau(w)=1\).
\end{proof}

La forma cuadrática \(N_\tau\) es positiva definida únicamente en el tipo elíptico.
Para \(\tau=0\), vale \(a^2\), y el elemento \(J_0\) es no nulo aunque su
cuadrado y su norma sean cero. Para \(\tau=-1\), vale \(a^2-b^2\); los
elementos no nulos \(I+J_{-1}\) e \(I-J_{-1}\) tienen norma cero y su
producto se anula. La norma algebraica es aquí una forma cuadrática
multiplicativa: no presupone una norma positiva de espacio vectorial.

Los tres tipos pueden realizarse fielmente mediante
\[
 aI+bJ_\tau\longmapsto
 \begin{pmatrix}a&-\tau b\\b&a\end{pmatrix}.
\]
El determinante de esta matriz es \(N_\tau(aI+bJ_\tau)\). Quedan
representadas así, con una misma expresión matricial, la estructura
compleja, la estructura de números duales y la estructura real escindida.
Cada una conserva su ley y sus divisores de cero. La forma de firma
\((1,1)\) del último caso tiene interpretación lorentziana en dimensión
\(1+1\); la identificación de una métrica física requiere además un espacio,
un campo y una regla de transporte específicos.

\subsubsection{Conservación multiplicativa y composición de evoluciones}
\label{trit-ampl-evoluciones}

La exponencial ya obtenida pertenece al conjunto de elementos de norma uno:
\[
 N_\tau\bigl(\exp(tJ_\tau)\bigr)=1.
\]
Esta conservación se mantiene al componer parámetros. En efecto, todos los
factores son funciones del mismo generador, y
\[
 \exp(tJ_\tau)\exp(sJ_\tau)=\exp((t+s)J_\tau).
\]
Al comparar las dos componentes se deducen las leyes de adición
\begin{align}
 C_\tau(t+s)&=C_\tau(t)C_\tau(s)-\tau S_\tau(t)S_\tau(s),\\
 S_\tau(t+s)&=S_\tau(t)C_\tau(s)+C_\tau(t)S_\tau(s).
\end{align}
La conjugación cambia \(t\) por \(-t\) y coincide con la inversión dentro
de este subgrupo. No cambia \(\tau\): invierte una evolución del régimen
elegido.

En el tipo elíptico, la forma conservada es \(a^2+b^2\), y la trayectoria
exponencial recorre el círculo de norma uno. En el hiperbólico, las
coordenadas propias son \(u=a+b\), \(v=a-b\), con \(uv=N_{-1}(z)\).
La exponencial da
\[
 (u,v)=(e^t,e^{-t}),\qquad uv=1.
\]
El crecimiento de una coordenada se acompaña de la contracción de la otra.
La conservación expresa una relación entre ambas coordenadas; no afirma que
cada una permanezca constante.

En el tipo parabólico,
\[
 (I+tJ_0)(I+sJ_0)=I+(t+s)J_0,\qquad
 (I+tJ_0)^{-1}=I-tJ_0.
\]
La evolución no se detiene porque \(J_0^2=0\). La nilpotencia elimina los
términos de orden al menos dos, pero conserva el término lineal que
transporta el parámetro. Son distintos el símbolo visible \(0\), el
elemento nulo del álgebra y el generador no nulo \(J_0\). Esta distinción
es necesaria para describir el umbral sin identificarlo con ausencia de
estado o de dinámica.

\begin{proposition}[Composición de coordenadas afines]
\label{trit-ampl-prop-pendientes}
En la coordenatización donde la componente escalar no se anula, represéntese una
dirección mediante \(I+uJ_\tau\). Si \(1-\tau uv\ne0\), la dirección de su
producto con \(I+vJ_\tau\) tiene coordenada
\[
 u\star_\tau v=\frac{u+v}{1-\tau uv}.
\]
\end{proposition}
\begin{proof}
Se tiene
\[
 (I+uJ_\tau)(I+vJ_\tau)
   =(1-\tau uv)I+(u+v)J_\tau.
\]
Dividir por la componente escalar da la expresión anunciada.
\end{proof}

La coordenatización afín conserva la razón entre componentes, no el factor escalar
dividido. Cuando \(1-\tau uv=0\), debe examinarse el producto antes de
cambiar de coordenatización: puede tener dirección definida con componente escalar
nula, o ser el elemento cero en presencia de divisores de cero. La ley de
composición sigue perteneciendo al álgebra completa; la singularidad de
una coordenada no equivale automáticamente a una singularidad del objeto.
Esta fórmula reaparecerá en la lectura regional elíptica. Su uso requiere
mantener el denominador y, cuando corresponda, el factor de normalización.

\subsubsection{Transporte del tipo y realizaciones posteriores}
\label{trit-ampl-puentes}

Las álgebras \(\mathbb A_{+1}\), \(\mathbb A_0\) y
\(\mathbb A_{-1}\) no son isomorfas como álgebras reales unitarias.
La primera es un cuerpo; la segunda contiene un nilpotente no nulo; la
tercera posee idempotentes no triviales y es isomorfa a
\(\mathbb R\times\mathbb R\). Por ello, el cambio de firma que organiza una
ruta no puede interpretarse como un isomorfismo indiferenciado entre las
tres realizaciones. El TPK conserva el índice de régimen, los dominios y el
operador efectivo que actúa en cada transición.

La aplicación electrónica desarrollada en el Artículo I toma el par de
proyecciones de las hojas de APP y construye sobre él un conmutador
normalizado cuyo cuadrado es \(-I\). El mismo bloque realiza una involución de
cuadrado \(I\) y dos direcciones nilpotentes. El presente artículo reutiliza
este resultado operatorio para tipar la estructura compleja de las dos-formas
y la respuesta de Cartan--Holst; la clasificación precedente establece por
qué las tres relaciones cuadráticas coexisten en una misma álgebra matricial
conservando sus dominios y sus acciones respectivas.

En la prolongación excepcional, la matriz de Paley reducida módulo \(3\)
satisface \(A_W^2=-I_6\). La relación dota al espacio ternario de una
estructura de dimensión \(3\) sobre \(\mathbb F_9\). Su vínculo con una raíz
operatorial real se expresa mediante la presentación integral
\[
 \mathbb Z[u]/(u^2+1).
\]
Los cambios de base a \(\mathbb F_3\) y a \(\mathbb R\) producen
representaciones de esa misma relación. Cada una conserva su característica,
su dimensión y su información de incidencia. El desarrollo de los
respectivos operadores establece el enlace, no la igualdad de sus
cardinales.

La conservación de norma uno describe, por consiguiente, las realizaciones
cuadráticas de las evoluciones. La conservación de cocientes, orientación y
trayectoria pertenece al estado aritmético transportado. Ambas propiedades
intervienen en la construcción, con funciones diferentes: una determina las
identidades de las representaciones; la otra permite componer, reconstruir y
prolongar las operaciones que las originan.
